% GENERATED FILE. Edit canonical lessons, then run npm run generate:books.
% canonical-source-hash: fnv1a64:1a9370fd1e64fb44
% canonical-lessons: BN-C19-je, BN-C19-jokhon, BN-C19-abar-bolben, BN-C19-dhire, BN-R19-lost-and-found

\chapter{When You Are Lost}
\label{ch:bn-when-you-are-lost}
\hlchaptermodality{30}

\begin{chapteropening}
\textbf{By the end of this chapter:} \emph{I can report what I think with \bn{যে}, put one event beside another in time with \bn{যখন} … \bn{তখন}, and ask somebody to say it again, slowly, when I have not followed.}

A learner who has stopped following says so and gets the sentence back: I do not understand, say it again, slowly — assembled out of \bn{না} from chapter twenty-seven and two words the book taught in chapters seven and eleven — and then reports an opinion with \bn{যে} and places it in time with \bn{যখন} … \bn{তখন}.
\end{chapteropening}

\section[je]{\bn{যে} (je) — that}

\label{lesson:BN-C19-je}



\begin{quote}
\textbf{\bn{ভাবা}}, \textquotedblleft{}to think\textquotedblright{}, has been in this book since chapter fourteen, and in all that time nobody has been able to say what they think.

\begin{itemize}
\raggedright
  \item \emph{Recall:} two items back, the word that heads a reason — \emph{kenonā}
\end{itemize}
\end{quote}

\subsection*{You'll want to know: \bn{যে}}
\begin{quote}
\textbf{\bn{আমি} \bn{ভাবি} \bn{যে} \bn{চা} \bn{ভালো}} — \emph{āmi bhābi je chā bhālo} — \textbf{I think that tea is good.}
\end{quote}

\begin{quote}
\textbf{\bn{আমি} \bn{জানি} \bn{যে} \bn{তুমি} \bn{ভালো}} — \emph{āmi jāni je tumi bhālo} — \textbf{I know that you are well.}
\end{quote}

\textbf{\bn{যে}} does one job: it takes a whole sentence and turns it into the object of a verb of thinking, knowing or saying. Everything before it is the frame; everything after it is the thought, unchanged, in the order it would have had standing alone.

\begin{grammarlens}[title={Grammar Lens: what English drops and Bengali keeps}]
English can leave \emph{that} out — \emph{I think tea is good} is ordinary English. Bengali speech drops \textbf{\bn{যে}} too, and Bengali writing keeps it. A learner should say it: the sentence is easier to hear coming when the joint is marked.

The word is the Sanskrit relative stem \emph{ya-}, an inheritance shared with Latin \emph{qui} and English \emph{who}. That is why the same syllable will later introduce a relative clause as well — one word, two jobs, and this lesson teaches the simpler one.
\end{grammarlens}

\subsection*{Your turn}
Say these aloud:

\begin{itemize}
\raggedright
  \item \emph{āmi bhābi je chā bhālo}
  \item \emph{āmi jāni je tumi bhālo}
  \item an opinion, then an objection to it — \emph{āmi bhābi je chā bhālo, kintu dudh nôy}
\end{itemize}

\begin{tcolorbox}[breakable,colback=teal!4,colframe=teal!35!black,title={Before you move on}]
What does \textbf{\bn{যে}} attach to the verb in front of it? (\textbf{A whole sentence}, as its object.) Does the sentence after it change shape? (\textbf{No} — it stands exactly as it would alone.) And which verb, taught in chapter fourteen, has been waiting for this word? (\textbf{\bn{ভাবা}}, to think.)

Source: \href{https://www.unicode.org/charts/PDF/U0980.pdf}{Unicode Bengali chart}.
\end{tcolorbox}
\section[jôkhon]{\bn{যখন} (jôkhon) — when}

\label{lesson:BN-C19-jokhon}



\begin{quote}
The lesson before this took the relative stem \textbf{\bn{য}} and hung a sentence off a verb with it. The same stem, with an ending, hangs a sentence off a moment.

\begin{itemize}
\raggedright
  \item \emph{Recall:} the complementiser from the lesson before — \emph{je}
\end{itemize}
\end{quote}

\subsection*{You'll want to know: \bn{যখন}}
\begin{quote}
\textbf{\bn{যখন} \bn{আমি} \bn{আসি}, \bn{তখন} \bn{চা} \bn{খাই}} — \emph{jôkhon āmi āshi, tôkhon chā khāi} — \textbf{when I come, then I drink tea.}
\end{quote}

Bengali does not leave the pair half-finished the way English does. \textbf{\bn{যখন}} opens the clause and \textbf{\bn{তখন}} answers it, and a fluent speaker says both. English would say \emph{when I come, I drink tea} and let the \emph{then} be understood; Bengali puts it in.

\begin{grammarlens}[title={Grammar Lens: one frame, three prefixes}]
\noindent
\begin{tabularx}{\linewidth}{@{}>{\raggedright\arraybackslash}X>{\raggedright\arraybackslash}X>{\raggedright\arraybackslash}X@{}}
\toprule
\textbf{word} & \textbf{prefix} & \textbf{job} \\
\midrule
\textbf{\bn{যখন}} & \textbf{\bn{য}-}, the relative & when … (opens a clause) \\
\textbf{\bn{তখন}} & \textbf{\bn{ত}-}, the demonstrative & then (answers it) \\
\textbf{\bn{কখন}} & \textbf{\bn{ক}-}, the interrogative & when? (asks) \\
\bottomrule
\end{tabularx}

One ending, three heads. The same three letters begin \textbf{\bn{যে}} (relative), \textbf{\bn{তুমি}} (demonstrative in origin) and \textbf{\bn{কি}} and \textbf{\bn{কেমন}} (interrogative) — which is to say the reader has been meeting this pattern since chapter three without being shown it.
\end{grammarlens}

\subsection*{Your turn}
Say these aloud:

\begin{itemize}
\raggedright
  \item \emph{jôkhon āmi āshi, tôkhon chā khāi}
  \item the pair on its own — \emph{jôkhon … tôkhon}
  \item a reason and a time in one breath — \emph{kenonā jôkhon āmi jāi …}
\end{itemize}

\begin{tcolorbox}[breakable,colback=teal!4,colframe=teal!35!black,title={Before you move on}]
Which word answers \textbf{\bn{যখন}}, and does a Bengali speaker leave it out? (\textbf{\bn{তখন}}, and no — both halves get said.) What do \textbf{\bn{যখন}}, \textbf{\bn{তখন}} and \textbf{\bn{কখন}} share? (\textbf{One ending}, with the relative, demonstrative and interrogative heads in front of it.) And which word from the previous lesson shares the first of those heads? (\textbf{\bn{যে}}.)

Source: \href{https://www.unicode.org/charts/PDF/U0980.pdf}{Unicode Bengali chart}.
\end{tcolorbox}
\section[ābār bôlben]{\bn{আবার} \bn{বলবেন} — say it again}

\label{lesson:BN-C19-abar-bolben}



\begin{quote}
\textbf{\bn{আবার}} was taught in chapter seven and \textbf{\bn{বলা}} in chapter eleven. Nineteen chapters later they have never stood next to each other, and the sentence they make is the one a beginner needs more than any other.

\begin{itemize}
\raggedright
  \item \emph{Recall:} two items back, the complementiser — \emph{je}
\end{itemize}
\end{quote}

\subsection*{You'll want to know: \bn{আবার} \bn{বলবেন}}
\begin{quote}
\textbf{\bn{আবার} \bn{বলবেন}} — \emph{ābār bôlben} — \textbf{please say it again.}
\end{quote}

\begin{quote}
\textbf{\bn{দয়া} \bn{করে} \bn{আবার} \bn{বলবেন}} — \emph{doyā kore ābār bôlben} — the same request, with \emph{please} in front of it.
\end{quote}

The ending is not new. Chapter twenty-two taught \emph{māf kôrben}, \textquotedblleft{}excuse me\textquotedblright{}, and named \textbf{-\bn{বেন}} as the polite ending that sounds like a future and works as a request. Put it on \textbf{\bn{বলা}} and you have asked somebody to repeat themselves without learning a single new word.

\begin{grammarlens}[title={Grammar Lens: the whole repair kit, in two sentences}]
A learner who is lost needs to say two things and no more: \emph{I do not understand} and \emph{say it again}. Chapter twenty-seven built the first out of \textbf{\bn{না}}. This lesson builds the second out of two words from chapters seven and eleven. Neither needed new vocabulary; both needed somebody to put the pieces together.

\textbf{\bn{দয়া} \bn{করে}} goes in front when the request should be softer, exactly as chapter twenty-two taught it.
\end{grammarlens}

\subsection*{Your turn}
Say these aloud:

\begin{itemize}
\raggedright
  \item \emph{ābār bôlben}
  \item \emph{doyā kore ābār bôlben}
  \item the pair a lost learner needs — \emph{āmi bujhi nā … ābār bôlben}
\end{itemize}

\begin{tcolorbox}[breakable,colback=teal!4,colframe=teal!35!black,title={Before you move on}]
Which ending makes this a polite request, and where was it taught? (\textbf{-\bn{বেন}}, in chapter twenty-two, on \emph{māf kôrben}.) How many new words does this sentence need? (\textbf{None} — both were already in the book.) And what goes in front to soften it? (\textbf{\bn{দয়া} \bn{করে}}.)

Source: \href{https://www.unicode.org/charts/PDF/U0980.pdf}{Unicode Bengali chart}.
\end{tcolorbox}
\section[dhīre]{\bn{ধীরে} (dhīre) — slowly}

\label{lesson:BN-C19-dhire}



\begin{quote}
Asking somebody to repeat themselves only helps if they repeat themselves at a speed you can follow.

\begin{itemize}
\raggedright
  \item \emph{Recall:} the request from the lesson before — \emph{ābār bôlben}
\end{itemize}
\end{quote}

\subsection*{You'll want to know: \bn{ধীরে}}
\begin{quote}
\textbf{\bn{ধীরে} \bn{বলবেন}} — \emph{dhīre bôlben} — \textbf{please speak slowly.}
\end{quote}

\begin{quote}
\textbf{\bn{ধীরে} \bn{ধীরে} \bn{বলবেন}} — \emph{dhīre dhīre bôlben} — the same thing, said the way Bengali actually says it.
\end{quote}

The doubling is not emphasis for a foreigner's benefit. Bengali turns an adverb up by repeating it, and \textbf{\bn{ধীরে} \bn{ধীরে}} is the ordinary form of this word in speech; the single \textbf{\bn{ধীরে}} is what a dictionary lists.

\begin{cousinweb}
\textbf{\bn{ধীরে}} is Sanskrit \emph{dhīra}, which meant \textbf{steady, composed, deliberate} long before it meant \emph{slow}. The word is about the manner of a person, and only then about the speed of a thing — which is why asking somebody to speak \emph{dhīre} is a gentler request than asking them to speak slowly in English.

Its second letter carries the long \emph{i}-sign chapter twenty-four taught on \textbf{\bn{নীল}}. Four pieces, all of them already read.
\end{cousinweb}

\subsection*{Your turn}
Say these aloud:

\begin{itemize}
\raggedright
  \item \emph{dhīre bôlben}
  \item the doubled form — \emph{dhīre dhīre bôlben}
  \item the whole repair kit — \emph{āmi bujhi nā … ābār bôlben … dhīre dhīre}
\end{itemize}

\begin{tcolorbox}[breakable,colback=teal!4,colframe=teal!35!black,title={Before you move on}]
What did \textbf{\bn{ধীরে}} mean in Sanskrit before it meant \emph{slowly}? (\textbf{Steady, composed} — a word about a person.) Why do Bengali speakers say it twice? (\textbf{Doubling is how the language intensifies an adverb}.) And which sign sits on its second letter? (\textbf{The long \emph{i}}, from \textbf{\bn{নীল}}.)

Source: \href{https://www.unicode.org/charts/PDF/U0980.pdf}{Unicode Bengali chart}.
\end{tcolorbox}
\section[je · jôkhon · ābār bôlben]{When you are lost}

\label{lesson:BN-R19-lost-and-found}



\begin{quote}
Four items landed in this chapter. Two of them join clauses and two of them rescue a conversation. Name all four.
\end{quote}

\subsection*{Your turn}
\begin{itemize}
\raggedright
  \item \emph{Say it:} an opinion — \emph{āmi bhābi je chā bhālo}
  \item \emph{Say it:} a time — \emph{jôkhon āmi āshi, tôkhon chā khāi}
  \item \emph{Say it:} the repair kit in order — \emph{āmi bujhi nā … ābār bôlben … dhīre dhīre}
  \item \emph{Recall:} eight lessons back, the word that heads an objection — \emph{kintu}
  \item \emph{Recall:} fourteen lessons back, the joiner that adds — \emph{ār}
\end{itemize}

\begin{tcolorbox}[breakable,colback=teal!4,colframe=teal!35!black,title={Before you move on}]
Four chapters ago this book could not deny a sentence, join two nouns, give a reason or ask for a repeat. Say the sentence that agrees and then objects. (\textbf{\bn{আমি} \bn{একমত}, \bn{কিন্তু} …}) Say the one that reports a thought. (\textbf{\bn{আমি} \bn{ভাবি} \bn{যে} …}) And say the two sentences a lost learner needs. (\emph{\textbf{āmi bujhi nā}} and \textbf{\bn{আবার} \bn{বলবেন}}.)
\end{tcolorbox}
